\section{\iflanguage{english}{Conclusion}{Fazit}}\label{sec:conclusion}

This project investigated the mechanisms behind plasticity loss in neural networks, with a particular focus on deep reinforcement learning. The main objective was to understand why networks that have already undergone training often become less capable of adapting to new learning signals. To address this question, we studied the recent literature on plasticity loss, capacity loss, warm-starting, and value-function learning, and reproduced several central experiments from \citep{lyle2023understanding} to verify and get the understanding of what the authors have conveyed in this seminal work of theirs.

After reproducing these experiments in this project, we can support the view that plasticity loss is closely connected to the geometry of the loss landscape. In particular, the optimizer instability experiment showed that abrupt changes in the learning objective can significantly disturb the gradient structure and cause many hidden ReLU units to become inactive. This prevents the network from effectively using its representational capacity and can strongly reduce future learning ability. The Hessian and gradient covariance experiments further showed that training can move the network into regions of parameter space with sharper curvature and stronger gradient interference. These observations suggest that plasticity loss is not only a representational problem, but also an optimization problem.

Based on the reproduced experiments and the methods studied in \citep{lyle2023plasticity}, the most effective interventions were those that made loss curvature easier to optimize. In particular, methods such as layer normalization and categorical output representations were more successful than methods that only perturbed or regularized the parameters. This supports the hypothesis that preserving plasticity requires controlling the sharpness and conditioning of the loss landscape.

Layer normalization is especially relevant in this context because it changes the parameterization of the network in a way that can stabilize optimization. By normalizing intermediate activations, layer normalization can reduce sensitivity to scale changes and help prevent the network from entering extremely sharp regions of the loss landscape. In experiments of \citep{lyle2023understanding}, the addition of layer normalization reduced plasticity loss in several architectures. This suggests that normalization methods may help preserve plasticity by producing smoother and better-conditioned optimization dynamics.

The results are also consistent with the conclusions of \citep{farebrother2024stop}, who argue that value functions in deep reinforcement learning should not necessarily be trained as standard regression problems. In conventional value-based RL, the network is trained to regress directly to bootstrapped target values. This can be difficult because the targets are noisy, non-stationary, and unbounded. Regression objectives may therefore push the parameters into regions of the loss landscape from which later adaptation becomes difficult. In contrast, categorical methods such as two-hot encoding or HL-Gauss transform value learning into a classification-style problem. Instead of directly predicting a scalar value, the network predicts a distribution over a fixed support. This can make the learning problem better conditioned and more stable.

Further works can study methods to use to reduce the sharpness or smoothen the navigation of network in function space so that plasticity remains intact.